% from main.tex
%%%%%%%%%%  scan idx 104  |  folio 91  %%%%%%%%%%
hypothèse de finitude sur \(\Ho^n(G)\) est conservée; une nouvelle correction locale en degré \(n\) permet alors d'obtenir \(\Ho^n(u)\) surjectif.

La suite exacte de cohomologie
\[
  \longrightarrow \Ho^n(G)\longrightarrow
  \Ho^{n+1}(E)\longrightarrow
  \Ho^{n+1}(F)\longrightarrow
\]
montre que le noyau de \(\Ho^{n+1}(u)\) est de \(\mathcal C_0\)-type fini. En d'autres termes, l'ensemble des \(X\) au-dessus de \(S\) tels qu'il existe une suite exacte
\[
  M\longrightarrow \Ho^{n+1}(E_X)
  \xrightarrow{\Ho^{n+1}(u)}
  \Ho^{n+1}(F_X),
\]
avec \(M\in\Ob\mathcal C_{0,X}\), est un raffinement de \(S\). Comme \(\mathcal C_0\) est localement quasi-relevable dans \(\mathcal C\), on peut d'ailleurs imposer que la flèche \(M\longrightarrow \Ho^{n+1}(E_X)\) se factorise à travers \(Z^{n+1}(E_X)\) en
\[
  M\xrightarrow{r}Z^{n+1}(E_X)\longrightarrow \Ho^{n+1}(E_X).
\]
La flèche composée
\[
  M\xrightarrow{u_X^{n+1}r}Z^{n+1}(F_X)\longrightarrow \Ho^{n+1}(F_X)
\]
étant nulle par construction, \(u_X^{n+1}r\) se factorise à travers \(B^{n+1}(F_X)=\Im\!\left(F_X^n\xrightarrow{d^n}F_X^{n+1}\right)\). Appliquant à nouveau la quasi-relevabilité, on peut imposer que \(u_X^{n+1}r\) se factorise à travers \(F_X^n\) en
\[
  M\xrightarrow{s}F_X^n\longrightarrow B^{n+1}(F_X)\longrightarrow Z^{n+1}(F_X).
\]
On vérifie alors trivialement que \(u'\), défini comme en (1.4.1), est un morphisme de complexes satisfaisant à (i).

Soit \(G'\) le cône de \(u':E'\to F_X\). Montrons que \(\Ho^n(G')\) est de \(\mathcal C_0\)-type fini. On a, par construction, une suite exacte
\[
  0\longrightarrow E_X\xrightarrow{v}E'\longrightarrow M[-n]\longrightarrow 0,
\]
